%%%%%%%%%%%%
%
% $Autor: Wings $
% $Datum: 2019-03-05 08:03:15Z $
% $Pfad: Maintenance.tex $
% $Version: 4250 $
% !TeX spellcheck = en_GB/de_DE
% !TeX encoding = utf8
% !TeX root = manual 
% !TeX TXS-program:bibliography = txs:///biber
%
%%%%%%%%%%%%

\chapter{Wartung}

\section{Einstellen und Austauschen von Komponenten} 
\label{sec:einstellen}

\begin{itemize}
	
	\item \textbf{Vorspannung ändern} \\
	Zur Erhöhung der Vorspannung die Einstellmuttern an den Seiten des Tischförderbands im Uhrzeigersinn drehen, bis die gewünschte Vorspannung erreicht ist. Um die Vorspannung zu verringern, die Einstellschrauben gegen den Uhrzeigersinn drehen.\\
	Achtung: Bei zu hoher Vorspannung kann es zu Beschädigungen der Antriebskomponenten kommen. Ist die Vorspannung zu niedrig, kann das Transportband durchrutschen oder von den Laufrollen gleiten.
	
	\item \textbf{Spur einstellen} \\
	Um die Spur und damit den Geradeauslauf des Transportbandes einzustellen, muss das Tischförderband aktiviert werden. Nun wird die maximale Bandgeschwindigkeit in eine beliebige Richtung eingestellt. Bei korrekter Vorspannung kann nun eine der beiden Einstellschrauben verstellt werden, um den Geradeauslauf zu beeinflussen. Es wird per Sichtprüfung die Mittigkeit des Transportbands im Lauf durch Drehen der Einstellschraube angepasst. Die Drehrichtung ergibt sich aus der Reaktion des Transportbandes.
	
	\item \textbf{Transportband wechseln} \\
	Um das Transportband zu wechseln, die Vorspannung verringern und das Transportband seitlich entnehmen. Vor Einbau des Neuteils die Einstellschrauben noch etwa 2 Umdrehungen weiter lösen, dann das neue Transportband auflegen und die Vorspannung wiederherstellen und justieren.
	
\end{itemize}



\section{Alle 3 Monate}

\begin{itemize}
	
	\item \textbf{Vorspannung kontrollieren} \\
	Vorspannung des Transportbandes prüfen. Wenn sich das Band mehr als 1 cm abheben lässt, bitte die Vorspannung des Transportbandes erhöhen (siehe.~\ref{sec:einstellen}).
	
	\item \textbf{Verschleißprüfung} \\
	Sichtprüfung des Transportbands auf Schnitte oder Risse. Wenn Beschädigungen vorhanden sind, ist das Transportband zu erneuern.
	
	\item \textbf{Geräuschprüfung} \\
	Laufruhe des Transportbands sowie Geräuschentwicklung der Komponenten testen. Hierfür das Tischförderband bei unterschiedlichen Bandgeschwindigkeiten betreiben. Die Geräuschentwicklung nimmt im Betrieb mit erhöhter Bandgeschwindigkeit gleichmäßig zu. Bei erhöhter Lautstärke Vorspannung prüfen und bei Bedarf verringern (siehe.~\ref{sec:einstellen}). Wenn das Problem weiterhin auftritt, Lagerung und Antrieb soweit möglich kontrollieren.
	
	\item \textbf{Prüfung der Spurhaltigkeit} \\
	Zum Prüfen das Tischförderband mit maximaler Bandgeschwindigkeit und ohne Last aktivieren. Dann einseitig die Einstellschrauben nutzen, um das Transportband mittig auf den Laufrollen auszurichten (siehe.~\ref{sec:einstellen}).
	
	\item \textbf{Prüfung der Sensorik} \\
	Während aktivem Transportbandlauf ohne Gegenstände auf dem Transportband die Lichtschranke durch einen beliebigen Gegenstand aktivieren. Die Reaktion des Tischförderbands muss sofort erfolgen. Wenn dies nicht der Fall ist, Neustart durchführen. Bei bleibendem Fehler Software und Elektronik prüfen lassen.
	
\end{itemize}



\section{Halbjährlich}

\begin{itemize}
	
	\item \textbf{Prüfung der Laufruhe} \\
	Prüfung des Bandes bei 50 \% und später 100 \% Bandgeschwindigkeit auf Laufruhe. Zunächst ohne Gegenstände, dann mit Gegenständen als zusätzliche Belastung. Wenn unruhiger Lauf oder Vibrationen im Betrieb festgestellt werden, sind die Transportbandvorspannung, die Laufrollen sowie die Lagerung zu kontrollieren. Wenn diese trotz unruhigem Lauf fehlerfrei sind, Transportband erneuern.

	
\end{itemize}
